\documentclass[11pt]{article}
\usepackage[a4paper,margin=25mm]{geometry}
\usepackage{fontspec}
\usepackage[hidelinks]{hyperref}
\setmainfont{Latin Modern Roman}
\pagestyle{plain}
\begin{document}
{\LARGE\bfseries ALIUS Bulletin - Guidelines for conversations - 2022\par}
\vspace{1em}
\section*{Alius Bulletin}
\section*{*}
\section*{Guidelines For Conversation}
If you want to conduct an open conversation for ALIUS Bulletin, please read the following guidelines carefully.\par

Conversation format\par

Conversation will be initially recorded, aiming at increasing the flow in the exchange of ideas and collective reflection. This will provide the interviewers and the interviewees to conduct a conversation under naturalistic conditions (which does not have to be publicly released in this format). The conversation will then be fully transcripted by an automatic algorithm; and transcripts will finally be edited by both authors and editors. All interviews are formatted and integrated at the end of the year in the annual Bulletin.\par

Duration of the conversation\par

The length of the recorded conversation is ideally 30 - 60 minutes in length. After transcription and editing, an abstract, keywords, highlights and bibliography will be added.\par

Content of the interview\par

You are not expected to carry out the interview in a journalistic style – asking short and non-technical questions. If needed, do not hesitate to ask long and detailed questions. Bear in mind that interviewees might take many things for granted. Therefore, your job, as the interviewer, is to make sure that your questions summarize in an understandable manner key concepts and facts that will facilitate the reading for non-specialists. So, when deemed relevant, do not hesitate to introduce key concepts and key data to contextualize a bit the questions you are asking. That being said, if you want to ask short questions, feel free to do so: ask long questions only when the intelligibility of what you are saying requires it.\par

Generally speaking, the interviews are intended for scholars. Therefore, you are not expected to ask general audience questions as a science journalist would do. However, each interview should be easily understandable to any scholar. For example, an anthropologist reading an interview conducted with a neuroscientist should easily understand most of the discussion, and conversely, a neuroscientist reading an interview conducted with an anthropologist should easily understand most of the discussion.\par

We invite you to have a look at the fifth issue of the Bulletin : a conversation with Olson \& Yaden hosted by George Fejer, to have a concrete example of how questions can be phrased. However, this model is just two illustrations of what can be done. Feel free to differ from this model: as long as you broadly follow the above recommendations, it will be perfect.\par

Modus operandi and requirements\par

1. Ask one of the editors of the Bulletin whether the researcher you want to interview fits well with the Bulletin’s editorial line.\par

2. Ask the researcher you want to interview whether they agree in principle to be interviewed (when you do so, let them know about the general concept of the Bulletin, about the other researchers interviewed in previous issues of the Bulletin, specify that the interview is done in a written way, describe the modus operandi, the schedule, etc.).\par

3. You can now write up your questions. Feel free to discuss them with one of the editors if needed, but this is not required.\par

4. Set up a date to have a recorded conversation with interviewee (e.g., via Jitsi or OBS, see here for instructions). Explain that the recording will be transcribed into a written format and that nothing be published without consenting to the final edit. Ask the interviewee if they would like to have the questions in advance.\par

4. Engage in a conversation with the interviewee. Include an informal opening where you ask the interviewee to say something about their background and/or personality. Ask them more specific questions about their work, and include follow-up questions whenever interesting themes rise up.

5. Once the conversation is over and recorded, use an automatic algorithm to fully transcript the interview (or ask the editors for assistance). Upload the recording to the ALIUS Google drive (link provided by the editors).\par

6. Once you have the transcript of the interview, you will next have to:\par

If necessary, compress questions and answers to reduce repetition, overlap and redundancy.\par

Edit the transcript to proofread, edit and organize it in a coherent piece.\par

Give a title to the interview (it is better to discuss this directly with the interviewee as you will need his/her final agreement for the title and the content of the interview).\par

Write a short abstract (<200 words) and up to 5 keywords for referencing the interview (you will need the final agreement of the interviewee for these).\par

If the interviewee discusses books or articles in the interview, it is your job to insert a bibliography at the end of the interview where all the books and articles mentioned by you (in your questions) or the interviewee (in his responses) are referenced. For the citation style that you are expected to use, see step 7 below.\par

Select 2 to 6 excerpts (depending on the length of the interview) highlighting some key points made by the interviewee and nicely epitomizing his/her ideas. These passages must be no longer than one or two sentences (to have a clearer idea as to how to choose these passages, see highlighted excerpts in the first issue of the Bulletin). Insert the 2, 3, 4 or 5, 6 selected excerpts at the very end of the text.\par

7. When the file is ready – when it includes the interview (questions and responses), a title, a bibliography and a few selected excerpts, send the whole file to the author for approval and bring more changes if needed.\par

8. Once interviewee and interviewers agreed on the final version, send the file to the editors. Note that the text in the file must comply with the following requirements:\par

Use Times New Roman, size 12.\par

Citation style for the bibliography is the 6th edition of APA (if you use Zotero, this citation style can be downloaded at the following link: https://www.zotero.org/styles?q=id\%3Aapa).\par

The Bulletin is written in American English and not in British English. If the text that you or the interviewee wrote is in British English, it will have to be turned into American English (e.g.: programme => program; realised => realized; etc.). This can be easily done using Word Autocorrect. If yourself or the interviewee uses British lexicon or phrasings, that’s fine: you can keep it as it is. But typographic details must be changed to fulfill the rules of American English. In addition to the “s” transformed into “z”, keep in mind that in American English, abbreviations such as “e.g.” and “i.e.” are followed by a comma (for example: “i.e., the mind is…” and not “i.e. the mind is…”).\par

Use the following quotation marks: “” (do not use «» or ‘’). If there is a quotation within a quotation, use the following quotation marks: ‘’. For example: “I know that the meaning he gives to ‘consciousness’ differs from…”\par

Make sure that there is no space before punctuation marks (; : , . ! ?). For example it should be: …because of two reasons: first, it doesn’t… and not: …because of two reasons : first, it doesn’t…\par

Dashes used for punctuation must be “em dashes” and not “en dashes” and they must not be separated by spaces (see: https://en.wikipedia.org/wiki/Dash). For example, it should be: …the concept of consciousness—thus defined—proves difficult… and not: …the concept of consciousness – thus defined – proves difficult…
(We cannot really expect interviewees to comply with all these requirements. So each interviewer will be in charge of making sure that the interviewee’s responses do comply with these requirements.)\par

Before sending the file to the editors, take time to carefully go through the interview to make sure there is no typo.\par

9. Once they receive the text of the interview, the editors will forward it to a native English-speaking proofreader who will take care of checking grammar and linguistic correctness. Note that the job of proofreaders is not to correct typos. As indicated above, typos must be corrected before sending the file to the editors.\par

10. The editors will next send you the feedback provided by the proofreaders. You will be expected to integrate to the interview the changes suggested by the proofreaders.\par

11. Once all changes have been integrated, read the whole interview again very carefully to make sure that there is no typo left.\par

12. When the final version is ready, send it to the interviewee to get his/her final agreement for publication.\par

13. Once you have the interviewee’s agreement for publication, send the final version of the interview to one of the editors. The interview will then be ready to be published!\par

(14.) If the interviewee doesn’t agree to publish the interview as it is and wants some changes to be made, make these changes and then send the revised version of the piece to the editors; they will forward it once more to the proofreaders to make sure there are no grammatical/linguistic problems. Once proofreaders have gone through the interview, we will proceed as described in step 9 above.\par

For any questions, feel free to ask editors (here below)\par

ALIUS contact email: aliusresearchgroup@gmail.com\par

Editorial team\par

Daniel Friedman (danielarifriedman@gmail.com); Matthieu Koroma (matthieu.koroma@gmail.com) ; George Fejer (gyorgyf@live.com)\par

\end{document}
